\section{R ile basit doğrusal regresyon}
\label{sec:r-regression}

\textbf{Soru.} Uyku süresinden kaygı puanını yordayan doğrusal modeli kurun;
katsayı aralıklarını, yedi saat için ortalama güven aralığını, bireysel tahmin
aralığını ve temel tanı grafiklerini üretin.

\begin{lstlisting}[style=prismR]
veri <- data.frame(
  uyku = c(4.5,5.0,5.2,5.5,5.8,6.0,6.1,6.3,6.5,6.7,
           6.8,7.0,7.1,7.3,7.5,7.6,7.8,8.0,8.2,8.5),
  kaygi = c(68,52,75,60,66,48,70,55,64,58,
            45,62,53,46,60,43,55,49,57,40)
)
model <- lm(kaygi ~ uyku, data = veri)
summary(model)
confint(model)

yeni <- data.frame(uyku = 7)
predict(model, newdata = yeni, interval = "confidence")
predict(model, newdata = yeni, interval = "prediction")

par(mfrow = c(2, 2))
plot(model)
par(mfrow = c(1, 1))
influence.measures(model)
\end{lstlisting}

\begin{outputguidebox}
\textbf{Çözüm ve kontrol değerleri.} Model
$\widehat{\text{Kaygı}}=89{,}5403-4{,}9836\,\text{Uyku}$ biçimindedir.
Eğim standart hatası 1,6290, $t(18)=-3{,}061$ ve $p=0{,}00675$'tir.
$R^2=0{,}3421$ olur. Yedi saat için nokta tahmini yaklaşık 54,66'dır;
ortalama ve bireysel aralıklar ayrı \texttt{predict} çağrılarından okunur.
\end{outputguidebox}

\textbf{Yorum.} 
\texttt{interval="confidence"} aynı uyku süresindeki ortalama kaygı için,
\texttt{interval="prediction"} yeni bir öğrencinin puanı için aralık verir.
İkinci aralık bireysel değişkenliği içerdiğinden daha geniştir. Model tanıları
eğim anlamlı bulundu diye atlanmamalıdır.

\textbf{Pekiştirme ve yanıt.} Veri 4,5--8,5 saat aralığındadır. On iki saat
için öngörü hesaplanabilse bile güçlü bir dışa taşırmadır ve bilimsel olarak
ayrıca gerekçelendirilmelidir.